GoBIS: An integrated framework to analyse the goal
and business process perspectives in information systems
Marcela Ruiz a,n
, Dolors Costal b
, Sergio España a
, Xavier Franch b
, Óscar Pastor a
a
PROS Research Centre, Universitat Politècnica de València, Camino de Vera s/n, Edificio 1F, 46022 Valencia, Spain
b GESSI Research Group, Universitat Politècnica de Catalunya (UPC), Calle Jordi Girona 1-3, Campus Nord, Edifici Omega,
08034 Barcelona, Spain
article info
Article history:
Received 3 October 2014
Received in revised form
21 February 2015
Accepted 30 March 2015
Available online 30 April 2015
Keywords:
Modelling language
Requirements engineering
Goal modelling
Business process modelling
Ontological analysis
Metamodel integration
Performance and usability analysis
i
n
iStar
Communication Analysis
abstract
Context: Organisational reengineering, continuous process improvement, alignment
among complementary analysis perspectives, and information traceability are some
current motivations to promote investment and scientific effort for integrating goal and
business process perspectives. Providing support to integrate information systems
analysis becomes a challenge in this complex setting.
Objective: The GoBIS framework integrates two goal and business process modelling
approaches: i
n
(a goal-oriented modelling method) and Communication Analysis (a
communication-oriented business process modelling method).
Method: In this paper, we describe the methodological integration of both methods with
the aim of fulfilling several criteria: i) to rely on appropriate theories; ii) to provide
abstract and concrete syntaxes; iii) to provide scenarios of application; iv) to develop tool
support; v) to provide demonstrable benefits to potential adopters.
Results: We provide guidelines for using the two modelling methods in a top-down
analysis scenario. The guidelines are validated by means of a comparative experiment and
a focus-group session with students.
Conclusions: From a practitioner viewpoint (modeller and/or analyst), the guidelines
facilitate the traceability between goal and business process models, the experimental
results highlight the benefits of GoBIS in performance and usability perceptions, and
demonstrate an improvement on the completeness of the latter having an impact on
efficiency. From a researcher perspective, the validation has produced useful feedback for
future research.
& 2015 Elsevier Ltd. All rights reserved.
1. Introduction
Organisations are aware of the importance of evolving to
keep pace with changes in the market, technology, environment, law, etc. [1]. As a result, continuous improvement and
reengineering have become common practices in information
system (IS) engineering. Understanding organisations and
their needs for change often requires several interrelated
perspectives [2,3]. The IS engineering community has contributed a number of modelling languages that are typically
oriented towards a specific perspective, requiring approaches
to their integration [4].
In this paper, we focus on extending a business process
perspective with intentional aspects of organisations. Business process modelling languages provide primitives to
specify work practice (i.e., activities, temporal constraints,
and resources). Despite the fact that processes are widely
Contents lists available at ScienceDirect
journal homepage: www.elsevier.com/locate/infosys
Information Systems
http://dx.doi.org/10.1016/j.is.2015.03.007
0306-4379/& 2015 Elsevier Ltd. All rights reserved.
n
Corresponding author. Tel.: þ34 96 387 70 00x83534.
E-mail addresses: lruiz@pros.upv.es (M. Ruiz),
dolors@essi.upc.edu (D. Costal), sergio.espana@pros.upv.es (S. España),
franch@essi.upc.edu (X. Franch), opastor@pros.upv.es (Ó. Pastor).
Information Systems 53 (2015) 330–345
accepted as a means to achieve organisational goals [5],
process models give little attention to the strategic dimension [6]. The analysis, prioritisation, and selection of organisational strategies are the scope of intentional modelling
languages, which focus on the business roles, their goals,
and their relationships.
Business processes and goals are intrinsically interdependent [7] and several works provide detailed arguments
in favour of combining both perspectives: (i) an integrated
approach allows understanding of the motivation for
processes [6]; (ii) in a non-integrated approach, goals
may be used to guide process design [8]; (iii) traceability
is enhanced, which is necessary for enterprise management [9] and facilitates the sustainability of organisations
[10]; (iv) integration also helps in identifying crossfunctional interdependencies during business change
management by supporting the identification of the goals
that motivate change and the analysis of their impact on
processes [8,10,11].
The GoBIS framework (Goal and Business Process
Perspectives for Information System Analysis) pursues this
aim by integrating a goal-oriented and a business processoriented modelling language. There are several criteria
that one would expect from modelling language integration. For the framework definition, we highlight the
following: (i) the languages to combine need to be
formally described; (ii) the integration itself should be
well founded in theory; (iii) it should clarify the scenarios
where the integrated approach can be applied and provide
some scenario-dependent guidelines; (iv) it should provide tool support; (v) empirical studies should demonstrate some benefits to potential adopters. These criteria
guided our research. A comparative review (see Section 2)
reveals that proposals with similar aims do not fulfil one or
several of the above-mentioned criteria, revealing that the
challenge remains open.
This paper presents our design science endeavour [12]
from the problem investigation to the solution design, the
implementation of a modelling tool, and the solution
evaluation. We have chosen to integrate the languages
proposed by i
n
[3] (a goal-oriented modelling method),
and Communication Analysis (CA) [13] (a communicationoriented business process modelling method). The reason
for choosing i
n
is its expressiveness to specify dependencies, with which we intend to trace strategic motivations
and processes. In the case of CA, we aim to get the most
out of the communicational techniques in order to analyse
business processes; its notation is not what is important,
but rather the underlying concepts and modularity guidelines. Moreover, some current business process modelling
suites use BPMN with a communicative approach. In
addition, the authors are competent in these languages
and are able to confront the challenge.
The contributions of the paper are the following:
 We present the iStar2ca guidelines v2.0 which are
intended for a top-down modelling scenario and are
an evolution of the iStar2ca guidelines v1.0 reported in
[14] (see the evolution time-line in Fig. 1). The guidelines design is a method engineering effort; throughout
the paper, we use the terminology introduced in [15].
 We report a comparative experiment and a focus-group
session, which was carried out with students and
whose feedback has been taken into account to produce the iStar2ca guidelines V2.0. The results demonstrate that the subjects' effectiveness (specifically,
business process model completeness) is greatly
improved by the use of the iStar2ca guidelines, without
compromising efficiency. Also the perceptions of the
usability of the guidelines are positive.
The paper is structured as follows. Section 2 compares
related works. In Section 3, we introduce and exemplify
the methods selected for integration. Section 4 describes
the research methodology and Section 5 details the process and proposal for integrating i
n
and CA. Section 6
presents guidelines for a top-down modelling scenario.
Section 7 describes the modelling tool and technical
support. Section 8 presents how the top-down scenario
guidelines were evaluated through a lab-demo, a comparative experiment, and a focus group session with
students. Finally, Sections 9 and 10 conclude with a
discussion and future lines of work.
2. Related work
The combination of different methods and models to
obtain a profund and comprehensive understanding of the
system to be produced has received much of attention in
the IS literature. On the one hand, some general frameworks have been proposed to reconcile different perspectives; for instance, Salay et al. [16] propose the concept of
macromodel for the development, comprehension, consistency management, and evolution of a collection of related
models. On the other hand, different approaches focus on
some of these perspectives (typically two) in a specific
domain of interest and provide a customised solution for
their particular case. This is the approach followed in the
GoBIS framework presented in this paper.
There are several related works that focus on the
integration of goal and business process perspectives in
the domain of business process management and maintenance. These approach goal-oriented business process
reengineering from diverse angles. We analyse these
approaches based on the criteria mentioned in Section 1,
which may take different values:
(a) Ontological foundation: none; conceptual framework
(concepts and their relationships are defined); based
Fig. 1. Evolution time-line for the iStar2ca guidelines.
M. Ruiz et al. / Information Systems 53 (2015) 330–345 331
امکان دانلود نسخه تمام متن مقالات انگلیسی
امکان دانلود نسخه ترجمه شده مقالات
پذیرش سفارش ترجمه تخصصی
امکان جستجو در آرشیو جامعی از صدها موضوع و هزاران مقاله
امکان دانلود رایگان ۲ صفحه اول هر مقاله
امکان پرداخت اینترنتی با کلیه کارت های عضو شتاب
دانلود فوری مقاله پس از پرداخت آنلاین
ISIArticles
مرجع مقالات تخصصی ایران
پشتیبانی کامل خرید با بهره مندی از سیستم هوشمند رهگیری سفارشات
☑
✓
دریافت فوری متن کامل مقاله